
\documentclass[a4paper]{article}
\usepackage{neilbib}
\input{neilutils}
\bibliographystyle{neilbib}
\title{COM6130 --- Tutorial 1\\Probability Theory and Poisson Distributions}
\author{Neil Lawrence}

\date{October 4, 2001}
\begin{document}

\maketitle
%
\begin{enumerate}

\item Leaves fall from a tree to the ground as a Poisson process at a rate of one every minute.

\begin{enumerate}
\item You arrive at 4:00 pm. During the next ten minutes, what is the probability that you see twelve leaves fall?

\item What is the probability that you see $n$ leaves fall where $n$ goes from 0 to 20? Use MATLAB! Plot the results on a graph.

\item Where does the curve peak? Compare the peak with the expected value, $\lambda t$.
\end{enumerate}

\item $X$ takes on four different values with differing probabilities, as shown in the following table.
\begin{center}
\begin{tabular}{c|cccc}
$X$ & 1 & 2 & 3 & 4 \\
\hline
$P(X)$& 0.05 & 0.2 & 0.3 & 0.45
\end{tabular}
\end{center}

\begin{enumerate}
\item What is the mean of $X$? 
\item What is the standard deviation of $X$?
\item What is the expectation of $\ln X$?
\end{enumerate}

\item Show that
\begin{equation}
\expectation{n^2} = \lambda t e^{-\lambda t} \sum_{n=1}^\infty \frac{n (\lambda t)^{(n-1)}}{(n-1)!},
\end{equation}
where $\expectation{n^2}$ represents the expectation of $n^2$ under a Poisson distribution.

Then by making the substitution $m = n - 1$ show that
\begin{equation}
\expectation{n^2} = (\lambda t)^2 + \lambda t.  
\end{equation}

What, therefore, is the variance of a Poisson distribution?

\emph{Hint: There are similarities with the proof of the mean under a Poisson distribution we went through in the lecture.}
\end{enumerate}

\end{document}
